\SetPageTheme{story}
\clearpage
\thispagestyle{empty}

\begin{center}
{\Huge\bfseries STRUCTURI DE CONTROL}\par
\vspace{0.6em}
{\large Programul nu merge mereu pe acelasi drum. Uneori asteapta, verifica, alege sau revine.}
\end{center}

\vspace{0.8em}

\begin{tcolorbox}[colback=white!8!cqink,colframe=cqcyan,coltitle=white,coltext=white,fonttitle=\bfseries,title=Modelul pe care il urmarim]
Un automat de bauturi este un exemplu foarte bun pentru structurile de control. El:
\begin{enumerate}
\item asteapta o moneda;
\item verifica daca suma este suficienta;
\item alege daca ofera produsul sau cere completare;
\item revine in starea de asteptare dupa ce actiunea s-a incheiat.
\end{enumerate}
\end{tcolorbox}

\vspace{0.6em}

\begin{center}
\begin{tikzpicture}[node distance=1.2cm, >=Latex, every node/.style={font=\small}]
  \node[draw, rounded corners, fill=white!10!cqink, text=white, minimum width=4.4cm, minimum height=1.0cm] (idle) {asteapta moneda};
  \node[draw, rounded corners, fill=white!10!cqink, text=white, minimum width=4.4cm, minimum height=1.0cm, below=of idle] (sum) {verifica suma};
  \node[draw, rounded corners, fill=white!10!cqink, text=white, minimum width=4.4cm, minimum height=1.0cm, below left=1.6cm and 2.3cm of sum] (more) {cere completare};
  \node[draw, rounded corners, fill=white!10!cqink, text=white, minimum width=4.4cm, minimum height=1.0cm, below right=1.6cm and 2.3cm of sum] (vend) {elibereaza produsul};
  \draw[->,cqcyan,very thick] (idle) -- node[right,text=white]{moneda introdusa} (sum);
  \draw[->,cqmagenta,very thick] (sum) -- node[left,text=white]{insuficient} (more);
  \draw[->,cqlime,very thick] (sum) -- node[right,text=white]{suficient} (vend);
\end{tikzpicture}
\end{center}

\clearpage
\thispagestyle{empty}

\begin{tikzpicture}[remember picture,overlay]
  \draw[cqcyan, line width=2.8pt, ->]
    ([xshift=1.6cm,yshift=-6.4cm]current page.north west)
      .. controls ([xshift=8.3cm,yshift=-5.2cm]current page.north west)
      and ([xshift=-4.0cm,yshift=-8.0cm]current page.north east)
      .. ([xshift=1.2cm,yshift=-11.2cm]current page.east);
  \draw[cqmagenta, line width=2.4pt, ->]
    ([xshift=1.3cm,yshift=-13.0cm]current page.north west)
      .. controls ([xshift=7.2cm,yshift=-14.6cm]current page.north west)
      and ([xshift=-4.0cm,yshift=-15.0cm]current page.north east)
      .. ([xshift=1.2cm,yshift=-17.0cm]current page.east);
\end{tikzpicture}

\begin{center}
{\Large\bfseries Automatul revine mereu la o stare controlata}
\end{center}

\vspace{0.8em}

\begin{tcolorbox}[colback=white!8!cqink,colframe=cqlime,coltitle=white,coltext=white,fonttitle=\bfseries,title=Cum traducem automatul in programare]
\begin{enumerate}
\item \textbf{Asteptarea} seamana cu o stare in care programul nu face nimic pana nu apare un eveniment.
\item \textbf{Verificarea} seamana cu o structura \texttt{daca}: suma este suficienta sau nu.
\item \textbf{Revenirea} seamana cu o repetitie: dupa fiecare actiune, automatul este gata sa primeasca o noua comanda.
\end{enumerate}

Exact aceasta combinatie de idee o vom regasi in jocuri: astepti o intrare, verifici o conditie, actualizezi starea, apoi revii la urmatorul pas.
\end{tcolorbox}

\vfill

\begin{checkpointbox}[Fir pedagogic]
In paginile urmatoare luam pe rand trei instrumente:
\begin{itemize}
\item \texttt{daca} pentru alegere;
\item \texttt{cat timp} pentru repetitie condusa de o stare;
\item \texttt{pentru} pentru parcurgere planificata.
\end{itemize}
\end{checkpointbox}

\clearpage

\SetPageTheme{simulation}
\PageHeader
  {Simularea deciziilor}
  {Vedem concret cum se schimba traseul cand modificam o conditie.}
  {Control}

\SimulatorCard{cq-i1-if-branch}{Ramificare in timp real}{Cum acelasi cod intra pe ramura A sau B in functie de o conditie booleana.}{Comuta conditia dintre adevarat/fals, apoi noteaza ce bloc a rulat si ce variabila a fost actualizata.}

\begin{guidebox}[Dupa scanare]
\begin{enumerate}
\item Ce variabila decide ramura?
\item Ce schimbare ramane vizibila la finalul fiecarui drum?
\item Ce regula ai adauga pentru a preveni un drum invalid?
\end{enumerate}
\AnswerSpace{6}
\end{guidebox}

\clearpage

\SetPageTheme{guided}
\PageHeader
  {Tipuri de structuri, comparate}
  {Comparatia clara te ajuta sa alegi structura potrivita pentru problema potrivita.}
  {Control}

\begin{conceptbox}[Comparatie rapida]
\begin{tabularx}{\linewidth}{>{\bfseries}lX}
\toprule
Structura & Cand o alegi \\
\midrule
\texttt{daca} & Cand exista o conditie care poate schimba drumul. \\
\texttt{cat timp} & Cand nu stii exact cate repetari vor fi, dar stii conditia de continuare. \\
\texttt{pentru} & Cand stii startul, limita si pasul de parcurgere. \\
\bottomrule
\end{tabularx}
\end{conceptbox}

\begin{semibox}[Exercitiu de selectie]
Pentru fiecare situatie, alege structura potrivita si explica in 1--2 randuri:
\begin{enumerate}
\item deschide usa doar daca ai cheie;
\item numara secundele pana la zero;
\item parcurge toate celulele de pe o linie de tabla.
\end{enumerate}
\AnswerSpace{7}
\end{semibox}

\clearpage
